\section{OSPF ECMP Failover}

\subsection{Objective}

The objective of this experiment was to validate OSPF Equal-Cost Multi-Path, or ECMP, failover.

The experiment checked whether a router with two equal-cost OSPF next-hops could continue forwarding traffic when one ECMP branch failed. The experiment also compared native OSPF ECMP behaviour with BFD-assisted OSPF ECMP behaviour.

\subsection{ECMP Route Under Test}

The ECMP route was observed on \texttt{hpe-r3}.

\begin{verbatim}
10.0.24.0/24
via 10.0.23.2 on eth3 weight 1
via 10.0.34.3 on eth0 weight 1
\end{verbatim}

The same route was also installed in the Linux kernel as a multipath route.

\begin{verbatim}
10.0.24.0/24 proto bird metric 32
nexthop via 10.0.23.2 dev eth3 weight 1
nexthop via 10.0.34.3 dev eth0 weight 1
\end{verbatim}

This proves that \texttt{hpe-r3} had two equal-cost OSPF forwarding branches:

\begin{verbatim}
hpe-r3 -> hpe-r2
hpe-r3 -> hpe-r4
\end{verbatim}

The failed ECMP next-hop was \texttt{10.0.23.2}. The surviving ECMP next-hop was \texttt{10.0.34.3}.

\subsection{Measurement Methods}

Multiple measurement methods were used because ECMP failover has both control-plane and data-plane behaviour.

The Linux kernel route monitor detected when the kernel forwarding route changed. Kernel route sampling confirmed when the failed ECMP next-hop was removed from the kernel route. The \texttt{ip route get} sampler showed the selected forwarding next-hop for the test destination. BIRD route-table sampling confirmed when BIRD removed the failed next-hop. ICMP ping measured packet loss during the failover window.

\subsection{OSPF-BFD Activation}

Initially, ECMP failover was measured using native OSPF behaviour. Then BFD was enabled on OSPF interfaces.

After enabling BFD, \texttt{hpe-r3} showed active BFD sessions to both ECMP neighbours:

\begin{verbatim}
10.0.23.2 eth3 Up
10.0.34.3 eth0 Up
\end{verbatim}

This confirms that BFD was active on the OSPF ECMP core links.

\subsection{Comparison Result}

\begin{table}[h]
\centering
\scriptsize
\begin{tabular}{|l|l|l|r|r|r|r|}
\hline
\textbf{Test} & \textbf{BFD} & \textbf{Ping} & \textbf{Kernel} & \textbf{Route Sample} & \textbf{BIRD} & \textbf{Loss} \\
\hline
Native OSPF ECMP & No & 0.02s & 840 ms & 820 ms & 819 ms & 16.00\% \\
\hline
OSPF ECMP with BFD, stress & Yes & 0.02s & 292 ms & 305 ms & 304 ms & 7.40\% \\
\hline
OSPF ECMP with BFD, normal & Yes & 1s & 517 ms & 542 ms & 540 ms & 1 packet \\
\hline
\end{tabular}
\caption{OSPF ECMP failover comparison}
\end{table}

\subsection{Interpretation}

The native OSPF ECMP baseline showed that the failed ECMP next-hop was removed, but convergence took longer.

After enabling BFD on OSPF interfaces, the ECMP failover improved significantly in the high-rate stress test.

\begin{verbatim}
Kernel monitor time improved from 840 ms to 292 ms.
BIRD route-table update improved from 819 ms to 304 ms.
Packet loss reduced from 16.00\% to 7.40\%.
\end{verbatim}

This shows that BFD helped OSPF detect the failed ECMP branch faster and remove the failed next-hop more quickly.

\subsection{Normal-Rate Traffic Result}

The normal-rate ping test used one ICMP packet per second.

\begin{verbatim}
20 packets transmitted
19 received
1 packet lost
5.00\% packet loss
\end{verbatim}

Although the percentage appears as 5\%, this represents only one lost packet during the failover window. This is a practical way to describe normal-rate traffic impact.

\subsection{Stress-Traffic Result}

The high-rate stress test used one packet every 0.02 seconds, or around 50 packets per second.

In the BFD-assisted stress test:

\begin{verbatim}
500 packets transmitted
463 received
37 packets lost
7.40\% packet loss
\end{verbatim}

At 50 packets per second, 37 lost packets corresponds to a short disruption window. This stress test made the failover window more visible.

\subsection{Route Behaviour}

Before failure, both ECMP next-hops were present:

\begin{verbatim}
nexthop via 10.0.23.2 dev eth3 weight 1
nexthop via 10.0.34.3 dev eth0 weight 1
\end{verbatim}

After the failure of the \texttt{10.0.23.2} branch, the kernel route changed to the surviving next-hop:

\begin{verbatim}
10.0.24.0/24 via 10.0.34.3 dev eth0 proto bird metric 32
\end{verbatim}

This proves that the failed ECMP branch was removed and forwarding continued through the surviving branch.

\subsection{Conclusion}

The OSPF ECMP failover experiment successfully demonstrated equal-cost multipath routing and failover.

Before failure, \texttt{hpe-r3} had two equal-cost next-hops toward \texttt{10.0.24.0/24}: \texttt{10.0.23.2} and \texttt{10.0.34.3}. When the \texttt{10.0.23.2} branch failed, OSPF removed the failed next-hop and continued forwarding using the surviving next-hop \texttt{10.0.34.3}.

The native OSPF ECMP baseline converged in 840 ms by kernel route monitor measurement. After BFD was enabled on OSPF interfaces, the stress-test convergence improved to 292 ms, and packet loss reduced from 16.00\% to 7.40\%.

Under normal-rate traffic, only one packet was lost during ECMP failover. This confirms that OSPF ECMP failover worked correctly, and BFD-assisted OSPF improved failure detection and reduced traffic disruption.
